% 2009 年山东卷文科第 18 题：直四棱柱 ABCD-A1B1C1D1，底面为等腰梯形。
% 来源 PDF：普通高考/2009/2009山东文.pdf，SHA256 1807901a21da57e3d45382294189b0c8e795722b6879842319ae75428639fad6
% 原页：第 2 页（A3 三栏）第 1 栏上部，印刷页 548。
% 重绘依据：底面 ABCD 按斜投影画成等腰梯形（AB 为下底、CD 为上底），A 左下、B 右下、C 右上、D 左上；
%   A1B1C1D1 为对应上底；E 在 AD 上、E1 在 AA1 上、F 在 AB 上；
%   被遮挡的棱 A1D1、D1C1、DD1 与辅助线画虚线。
\begin{tikzpicture}[scale=0.95, semithick, every node/.style={inner sep=1.3pt}]
  \coordinate (A) at (0,0);
  \coordinate (B) at (4.0,0);
  \coordinate (C) at (2.9,1.15);
  \coordinate (D) at (1.1,1.15);
  \coordinate (A1) at (0,2.2);
  \coordinate (B1) at (4.0,2.2);
  \coordinate (C1) at (2.9,3.35);
  \coordinate (D1) at (1.1,3.35);
  \coordinate (E) at ($(A)!0.5!(D)$);
  \coordinate (E1) at ($(A)!0.5!(A1)$);
  \coordinate (F) at ($(A)!0.5!(B)$);
  % 底面
  \draw (A) -- (B) -- (C) -- (D) -- cycle;
  % 上底（部分被遮挡）
  \draw (A1) -- (B1) -- (C1);
  \draw[dashed] (C1) -- (D1) -- (A1);
  % 侧棱
  \draw (A) -- (A1);
  \draw (B) -- (B1);
  \draw (C) -- (C1);
  \draw[dashed] (D) -- (D1);
  % 点
  \fill (E) circle (1.1pt);
  \fill (E1) circle (1.1pt);
  \fill (F) circle (1.1pt);
  \node[anchor=north west] at (A) {$A$};
  \node[anchor=north east] at (B) {$B$};
  \node[anchor=west] at (C) {$C$};
  \node[anchor=south east] at (D) {$D$};
  \node[anchor=east] at (A1) {$A_1$};
  \node[anchor=south east] at (B1) {$B_1$};
  \node[anchor=south west] at (C1) {$C_1$};
  \node[anchor=east] at (D1) {$D_1$};
  \node[anchor=north west] at (E) {$E$};
  \node[anchor=east] at (E1) {$E_1$};
  \node[anchor=north] at ([yshift=-2pt]F) {$F$};
\end{tikzpicture}
